\section{Strips}\label{sec:strip}

On two rows every rectangle is axis-aligned. Indeed, a tilted rectangle has no
horizontal side, so both sides at a lowest corner go strictly up, and the opposite
corner lies in a third row.

\begin{theorem}[Strip Theorem]\label{thm:strip}
Let $m\ge3$.
\begin{enumerate}[(a)]
\item A blackout of the $2\times m$ strip is valid if and only if at most one of its
columns is full.
\item We have $f(2,m)=m+1$, and the maximum blackouts are those with one full column
and one black point in every other column.
\item There are $a_2(m)=m\cdot2^{m-1}$ maximum blackouts.
\end{enumerate}
For $m=2$ there is only one rectangle, so every blackout is valid, $f(2,2)=4$, and
$a_2(2)=1$.
\end{theorem}

\begin{proof}
(a) Two distinct rectangles share at most one column. If they
share one, it cancels, and the difference set is the four points of the other two
columns $p$ and $q$, the \emph{pair constraint} on $\{p,q\}$. If they share
no column, the difference set is the eight points of four columns, and it contains a
pair constraint. Every pair constraint occurs, since $m\ge3$ gives a third column $z$,
and the rectangles on $\{z,p\}$ and $\{z,q\}$ have exactly this difference set. By
Corollary~\ref{cor:hitting}, $S$ is valid if and only if $K$ meets every two columns,
that is, if and only if at most one column has no kept point.

(b) By (a) a valid $S$ has $|K|\ge m-1$, so $|S|\le m+1$. Equality holds exactly when
one column is full and every other column has one kept point, hence one black point.
Such blackouts are valid by (a).

(c) There are $m$ choices for the full column and two in each other column.
\end{proof}

So row $2$ of \seqnum{A399580} is \seqnum{A001787}$(m)=m2^{m-1}$ from $m=3$ on.
Figure~\ref{fig:stripmax} shows one of the $192$ optima for $m=6$.

\begin{figure}[t]
\centering
\begin{tikzpicture}[x=0.700cm,y=0.700cm,baseline=(current bounding box.north)]
  \fill[keptcol] (3.6,-0.6) rectangle (4.4,-2.4);
  \draw[gridcol] (1,-1) grid (6,-2);
  \fill[black] (1,-1) circle (0.17);
  \filldraw[fill=white,draw=keptline,line width=0.6pt] (2,-1) circle (0.17);
  \filldraw[fill=white,draw=keptline,line width=0.6pt] (3,-1) circle (0.17);
  \fill[black] (4,-1) circle (0.17);
  \fill[black] (5,-1) circle (0.17);
  \filldraw[fill=white,draw=keptline,line width=0.6pt] (6,-1) circle (0.17);
  \filldraw[fill=white,draw=keptline,line width=0.6pt] (1,-2) circle (0.17);
  \fill[black] (2,-2) circle (0.17);
  \fill[black] (3,-2) circle (0.17);
  \fill[black] (4,-2) circle (0.17);
  \filldraw[fill=white,draw=keptline,line width=0.6pt] (5,-2) circle (0.17);
  \fill[black] (6,-2) circle (0.17);
  \node[lab] at (1,-0.35) {1};
  \node[lab] at (2,-0.35) {2};
  \node[lab] at (3,-0.35) {3};
  \node[lab] at (4,-0.35) {4};
  \node[lab] at (5,-0.35) {5};
  \node[lab] at (6,-0.35) {6};
  \node[lab] at (0.35,-1) {1};
  \node[lab] at (0.35,-2) {2};
\end{tikzpicture}

\caption{A maximum blackout of the $2\times6$ strip. Column $4$ (shaded) is full and
every other column has one black point, which gives $m+1=7$ points.}
\label{fig:stripmax}
\end{figure}

Counting by (a), the valid blackouts of the $2\times m$ strip have the size polynomial
$(1+2u)^m+m\,u^2(1+2u)^{m-1}$ for $m\ge3$.

\begin{remark}[orbits and nonattacking kings]\label{rem:kings}
Up to the row reflection and the column reversal, Burnside's lemma gives
$o_m=m\,2^{m-3}+[m\text{ odd}]\,2^{(m-3)/2}$ maximum blackouts of the $2\times m$ strip for
$m\ge3$. Only the identity and, for odd $m$ with the full column central, the reversal and
the half turn fix any. So $o_3,\ldots,o_8$ are $4,8,22,48,116,256$. These are the values
\seqnum{A322284}$(m-1)$, which count $m-1$ nonattacking kings on a $2\times(2m-2)$ board up
to symmetry. Put one king in each block of two board columns, the first $k$ on the left and
the rest on the right. The map makes column $k+1$ full and copies the rows of the kings into
the other columns. We do not write out the check that it is a bijection.


\end{remark}
